\section{Chapter 5: Nuclear Synthesis in Stars}

\subsection{Questions}
\begin{enumerate}[label=4.\arabic*]
    \item  In the simplest model of $\alpha$ decay, the $\alpha$ particle is pre-formed and trapped within the nucleus by a potential similar to a Coulomb barrier. The mean decay rate $\lambda$ is the frequency $\nu$ 
    (with which the $\alpha$ particle hits the barrier) times the probability of penetration of the barrier (quantum tunneling probability).
        \begin{enumerate}[label=(\alph*)]
            \item Write the approximate expression for the decay rate in terms of $\nu$, $E_G$, and $E$ (where $E$ is the energy released by $\alpha$ decay).
            \item Find the Gamow energy of the $\alpha$ decay of $^{235}\text{U}$.
            \item Find the Gamow energy of the $\alpha$ decay of $^{239}\text{Pu}$.
            \item Assuming that the frequency $\nu$ is approximately the same for both nuclei, find the ratio
           \[ \frac{\lambda(^{239}\text{Pu})}{\lambda(^{235}\text{U})} \]
           using $E(^{235}\text{U}) = 4.68 \textup{MeV}$ and $E(^{239}\text{Pu}) = 5.24 \textup{MeV}$. The lifetime of a radioactive species is inversely proportional to its decay rate. Given that the observed lifetime of $^{235}\text{U}$ is $\tau (^{235}\text{U}) = 7.1 \times 10^8 y$, estimate the lifetime of $^{239}\text{Pu}$. Compare your result with the experimental value
           \[ \tau (^{239}\text{Pu} ) = 2.41 \times 10^4 y \]
           and comment on the accuracy of this simple model.
        \end{enumerate}
    \item The \textbf{triple alpha process} is a set of nuclear reactions that occur in stars to produce carbon and other elements. Which of the following correctly describes the three steps of the triple alpha
          process, along with whether energy is absorbed or released during each step$?$
        \begin{enumerate}[label=(\alph*)]
            \item $^4\text{He} + ^4\text{He} \to ^8\text{Be}$
            \item $^8\text{Be} + ^4\text{He} \to ^{12}\text{C}^*$
            \item $^{12}\text{C}^* \to ^{12}\text{C} + \gamma$
        \end{enumerate}
\end{enumerate}

\subsection{Solutions}
\subsubsection*{Solution 4.1: $\alpha$ decay and Gamow energy}
\begin{enumerate}[label=(\alph*)]
    \item The decay rate $\lambda$ can be expressed as:
    \begin{equation}
    \lambda = \nu e^{-2\pi \sqrt{\frac{E_G}{E}}},
    \end{equation}
    where $\nu$ is the frequency of collisions with the barrier, $E_G$ is the Gamow energy, and $E$ is the energy released by $\alpha$ decay.
    
    \item The Gamow energy $E_G$ for the $\alpha$ decay of $^{235}\text{U}$ can be calculated using the formula:
    \begin{equation}
    E_G = (\pi \alpha Z_1 Z_2)^2 2 m_{r}c^2,
    \end{equation}
    where $Z_1 = 92$ (for uranium), $Z_2 = 2$ (for the alpha particle), and $m_r$ is the reduced mass of the system. The reduced mass can be approximated as:
    \begin{equation}
    m_r \approx \frac{m_{^{235}\text{U}} m_{\alpha}}{m_{^{235}\text{U}} + m_{\alpha}} \approx m_{\alpha},
    \end{equation}
    since the mass of uranium is much larger than that of the alpha particle. Substituting these values, we get:
    \begin{align*}
    E_G &\approx (\pi \times \frac{1}{137} \times 92 \times 2)^2 2 m_{\alpha} c^2 \\
        &\approx 25.7\,\text{MeV}.
    \end{align*}
    
    \item Similarly, for the $\alpha$ decay of $^{239}\text{Pu}$, we have $Z_1 = 94$ (for plutonium) and $Z_2 = 2$. The reduced mass can again be approximated as $m_{\alpha}$. Substituting these values, we get:
    \begin{align*}
    E_G &\approx (\pi \times \frac{1}{137} \times 94 \times 2)^2 2 m_{\alpha} c^2 \\
        &\approx 27.0\,\text{MeV}.
    \end{align*}
    
    \item The ratio of the decay rates can be calculated as follows:
    \begin{align*}
    \frac{\lambda(^{239}\text{Pu})}{\lambda(^{235}\text{U})} &= \frac{\nu e^{-2\pi \sqrt{\frac{E_G(^{239}\text{Pu})}{E(^{239}\text{Pu})}}}}{\nu e^{-2\pi \sqrt{\frac{E_G(^{235}\text{U})}{E(^{235}\text{U})}}}} \\
    &= e^{-2\pi \left( \sqrt{\frac{E_G(^{239}\text{Pu})}{E(^{239}\text{Pu})}} - \sqrt{\frac{E_G(^{235}\text{U})}{E(^{235}\text{U})}} \right)}.
    \end{align*}
    Substituting the values of $E_G$ and $E$, we get:
    \begin{align*}
    \frac{\lambda(^{239}\text{Pu})}{\lambda(^{235}\text{U})} &\approx e^{-2\pi \left( \sqrt{\frac{27.0}{5.24}} - \sqrt{\frac{25.7}{4.68}} \right)} \\
    &\approx e^{-2\pi \left( 2.27 - 2.35 \right)} \\
    &\approx e^{0.50} \approx 1.65.
    \end{align*}
    Therefore, the decay rate of $^{239}\text{Pu}$ is approximately 1.65 times that of $^{235}\text{U}$. Given that the lifetime $\tau$ is inversely proportional to the decay rate $\lambda$, we can estimate the lifetime of $^{239}\text{Pu}$ as follows:
    \begin{equation}
    \tau(^{239}\text{Pu}) \approx \frac{\tau(^{235}\text{U})}{1.65} \approx \frac{7.1 \times 10^8 y}{1.65} \approx 4.3 \times 10^8 y.
    \end{equation}
    Comparing this result with the experimental value of $\tau(^{239}\text{Pu}) = 2.41 \times 10^4 y$, we see that the simple model significantly overestimates the lifetime of $^{239}\text{Pu}$. This discrepancy suggests that the model does not capture all the complexities of $\alpha$ decay, such as nuclear structure effects and other factors that can influence the decay rate.
\end{enumerate}